\documentclass[12pt,a4paper]{article}
\usepackage[utf8]{inputenc}
\usepackage[T1]{fontenc}
\usepackage{amsmath,amssymb,amsfonts}
\usepackage{mathtools}
\usepackage{booktabs}
\usepackage{array}
\usepackage{geometry}
\usepackage{hyperref}
\usepackage{cite}
\geometry{margin=2.5cm}

\newcommand{\MeV}{\,\text{MeV}}
\newcommand{\GeV}{\,\text{GeV}}
\newcommand{\fpi}{f_\pi}
\newcommand{\qqbar}{\langle\bar{q}q\rangle}
\newcommand{\zch}{\mathfrak{z}}

\title{A Tale of Two Tuples: Seed Triples, Signed Charges, \\
and SUSY Breaking in the sBootstrap}

\author{A.~Rivero\thanks{Manuscript prepared with AI assistance (Claude/Anthropic). 
Continues the program of \cite{Rivero2005,Rivero2012}.}}
\date{March 2026}

\begin{document}
\maketitle

\begin{abstract}
We reformulate the sBootstrap in terms of signed mass charges 
$\zch = z_0 + z_i$ with $m = \zch^2$, and show that the Koide ratio 
$Q = 2/3$ is the algebraic consequence of two conditions: 
$\sum z_i = 0$ and $\langle z_i^2\rangle = z_0^2$ (perturbation 
energies average to the vacuum energy). Allowing negative charges, 
we find that pseudoscalar mesons satisfy a Koide triple 
$(-\pi, D_s, B)$ to 0.1\%, and identify \emph{seed triples} with 
one massless member---$(0, s, c)$ for quarks, $(0, \pi, D_s)$ for 
mesons---as the pre-breaking configurations from which the full 
spectrum emerges. These seed triples share the same analytical 
structure (charge ratio $2-\sqrt{3}$) and, strikingly, the meson 
seed $(0,\pi,D_s)$ has vacuum energy $v_0^2 \approx 351\MeV$, 
close to the lepton triple's $v_0^2 \approx 314\MeV$. A third 
tuple $(c,b,t)$ stands apart, having no meson counterpart because 
top diquarks do not hadronize. We discuss the diquark sector---which 
must fill the antisymmetric representation of $SU(5)$ flavor rather 
than the symmetric, requiring additional dynamical input---and outline 
how SUSY breaking proceeds through two mechanisms---F-term (chiral 
condensate, with Volkov--Akulov as its low-energy realization) and 
D-term (electromagnetic Fayet--Iliopoulos)---plus a separate 
heavy-sector origin for the $(c,b,t)$ triple.
\end{abstract}

\tableofcontents

%======================================================================
\section{Introduction}
\label{sec:intro}
%======================================================================

The sBootstrap \cite{Rivero2005,Rivero2012} proposes that the Standard 
Model's pseudoscalar mesons and leptons fill common $\mathcal{N}=1$ 
supermultiplets. Its foundational relation is
\begin{equation}
\label{eq:pimu}
m_\pi^2 - m_\mu^2 \approx f_\pi^2\,,
\end{equation}
identifying the pion decay constant $f_\pi \approx 92.2\MeV$ as the 
SUSY-breaking order parameter. Numerically, $m_\pi^2 - m_\mu^2 = 
8{,}316\MeV^2$ versus $f_\pi^2 = 8{,}501\MeV^2$, a $2.2\%$ 
discrepancy---comparable to the precision of the quark Koide triples 
discussed below. In this paper we develop two themes 
that change the structure of the program.

First, \textbf{mass is the square of a signed charge}: $m = \zch^2$ 
where $\zch = z_0 + z_i$ can be negative. The Koide ratio $Q = 2/3$ 
is then not numerology but the energy-balance condition 
$\langle z_i^2\rangle = z_0^2$. With signed charges, pseudoscalar 
mesons satisfy Koide, and the spectrum organizes into \emph{seed 
triples} with a massless member that grow into full triples upon 
symmetry breaking.

Second, \textbf{the full sBootstrap has three Koide tuples, not one}, 
and the breaking structure must generate all three: a light sector 
$(0,s,c) \to (-s,c,b)$ seen in both quarks and mesons, a lepton 
sector $(e\!\approx\!0, \mu, \tau)$, and a heavy sector $(c,b,t)$ 
that decouples from meson physics because the top quark does not 
hadronize. The diquark partners of quarks, needed to complete the 
SUSY multiplets, must fill the antisymmetric $\overline{10}$ of 
$SU(5)$ flavor rather than the symmetric 15, requiring extra 
dynamical input from the breaking sector.

%======================================================================
\section{The Koide Formula as Energy Balance}
\label{sec:koide_algebra}
%======================================================================

\subsection{Charges, not masses}

The fundamental variable is the signed mass charge $\zch = z_0 + z_i 
\in \mathbb{R}$, with the physical mass $m = \zch^2 \geq 0$.

\subsection{Derivation of $Q = 2/3$}

For three charges $\zch_k = z_0 + z_k$ with $z_0$ the mean 
($\sum z_k = 0$):
\begin{equation}
Q = \frac{\sum \zch_k^2}{(\sum \zch_k)^2} 
= \frac{3z_0^2 + \sum z_k^2}{9z_0^2}
= \frac{1}{3} + \frac{\langle z_k^2\rangle}{3z_0^2}\,.
\end{equation}
Therefore:
\begin{equation}
\boxed{Q = \frac{2}{3} \quad\Longleftrightarrow\quad 
\langle z_k^2\rangle = z_0^2\,.}
\end{equation}

The two conditions are:
\begin{enumerate}
\item \textbf{Zero-sum perturbations:} $z_1 + z_2 + z_3 = 0$ 
(definition of $z_0$ as the mean charge).
\item \textbf{Energy balance:} The mean energy of the perturbations 
equals the energy of the unperturbed charge: 
$\langle z_k^2\rangle = z_0^2$.
\end{enumerate}

Condition~2 is the physics. It says the splitting scale matches the 
base scale---the coefficient of variation equals unity: 
$\sigma(\zch)/\bar{\zch} = 1$. This is neither perturbative 
($\sigma \ll \bar{\zch}$, giving $Q \to 1/3$, all equal) nor 
dominant ($\sigma \gg \bar{\zch}$, giving $Q \to 1$, one mass 
dominates). The Koide value $2/3$ is the equipartition point.

\subsection{Why one mass can vanish}

If $\langle z_k^2\rangle = z_0^2$, a perturbation $z_k$ can reach 
$|z_k| \sim z_0$, making $\zch_k = z_0 + z_k \approx 0$. A vanishing 
mass is the \emph{natural extreme} of energy balance, not fine-tuning.

%======================================================================
\section{Seed Triples: The $(0, a, b)$ Structure}
\label{sec:seeds}
%======================================================================

\subsection{Analytical condition}

For a triple $(0, \zch_a, \zch_b)$ with one massless member, the 
Koide ratio becomes:
\begin{equation}
Q = \frac{\zch_a^2 + \zch_b^2}{(\zch_a + \zch_b)^2}\,.
\end{equation}
Setting $Q = 2/3$ and writing $r = \zch_a/\zch_b$ (with $r < 1$):
\begin{equation}
3(r^2 + 1) = 2(r + 1)^2 \quad\Longrightarrow\quad r^2 - 4r + 1 = 0\,,
\end{equation}
giving:
\begin{equation}
\boxed{r = 2 - \sqrt{3} \approx 0.26795\,.}
\end{equation}

Equivalently, in mass space: 
$m_{\text{light}}/m_{\text{heavy}} = (2 - \sqrt{3})^2 = 7 - 4\sqrt{3} 
\approx 0.07180$.

This is a \emph{universal} prediction: any Koide triple with one 
massless member has the same charge ratio between the two massive 
members, regardless of their absolute scale.

\subsection{The quark seed: $(0, s, c)$}

The strange and charm quark charges give:\footnote{Throughout 
this paper, quark masses are $\overline{\text{MS}}$ running masses 
evaluated at their own scale: $m_s(2\GeV) = 93.4\MeV$, 
$m_c(m_c) = 1.27\GeV$, $m_b(m_b) = 4.18\GeV$ (PDG 2024 central 
values). For the top quark, the pole mass $m_t = 172.76\GeV$ is 
used because it gives the best Koide fit ($Q = 0.6695$); the 
$\overline{\text{MS}}$ mass $\bar{m}_t(m_t) \approx 162.5\GeV$ 
gives $Q = 0.663$, significantly worse. This mixed convention 
reflects the fact that QCD running is large for light quarks and 
negligible for the top. Evaluating all masses at a common scale 
(e.g., $\mu = m_Z$) would change the light-quark Koide ratios 
by several percent, degrading the fit. Whether the ``natural'' 
scale for each quark is its own mass (as assumed here) or a 
universal scale is an open question tied to the seed-scale 
discussion in \S\ref{subsec:exactseed}.}
\begin{equation}
\frac{\sqrt{m_s}}{\sqrt{m_c}} = \frac{9.67}{35.64} = 0.2712\,, 
\qquad r_{\text{Koide}} = 0.2679\,, \qquad \text{deviation: } 1.2\%\,.
\end{equation}

This triple has $Q = 0.6644$, deviating from $2/3$ by only $0.35\%$. 
Its vacuum energy is $v_0^2 = (\zch_s + \zch_c)^2/9 \approx 228\MeV$.

A historical note: Harari, Haut, and Weyers \cite{Harari} observed 
a quark mass ratio pattern that, in the language of the present paper, 
corresponds to an approximate seed triple---$(0, d, s)$---at a time 
when the up quark mass was 
still thought to be possibly zero \cite{Harari}. With current masses, 
$\sqrt{m_d}/\sqrt{m_s} = 0.224$, deviating from $2 - \sqrt{3}$ by 
$16\%$---too large for a good Koide triple. But if $m_s$ were 
$\sim 150\MeV$ and $m_d \sim 7\MeV$ (values common in that era), the 
ratio improves considerably. The seed idea was already present, ahead 
of its time.

\subsection{The meson seed: $(0, \pi, D_s)$}

The pion and $D_s$ meson charges give:
\begin{equation}
\frac{\sqrt{m_\pi}}{\sqrt{m_{D_s}}} = \frac{11.81}{44.37} = 0.2663\,, 
\qquad r_{\text{Koide}} = 0.2679\,, \qquad \text{deviation: } 0.6\%\,.
\end{equation}

This gives $Q = 0.6679$, deviating by $0.18\%$. It is the 
\textbf{meson counterpart} of the quark seed $(0, s, c)$---both 
involve the strange-charm sector. Indeed, $D_s = c\bar{s}$ is the 
meson that directly contains the $s$ and $c$ quarks.

\subsection{The lepton triple as an ``almost-seed''}

The charged lepton triple $(e, \mu, \tau)$ has $\zch_e = 0.715\MeV^{1/2}$, 
which is small but nonzero. If $m_e$ were zero, we would need 
$\sqrt{m_\mu}/\sqrt{m_\tau} = 2 - \sqrt{3} = 0.2679$; the actual value 
is $0.2439$, deviating by $9\%$. So the leptons are \emph{not} an 
exact seed---the electron, while very light, is not massless. But the 
triple is clearly the same family: the near-critical Koide phase puts 
one member close to zero.

\subsection{Vacuum energy coincidence}

The $v_0^2$ values of the seed triples cluster remarkably:

\begin{center}
\begin{tabular}{lcc}
\toprule
Triple & $v_0^2$ (MeV) & $v_0$ ($\MeV^{1/2}$) \\
\midrule
$(e, \mu, \tau)$ & 314 & 17.72 \\
$(0, \pi, D_s)$ & 351 & 18.73 \\
$(0, s, c)$ & 228 & 15.10 \\
\bottomrule
\end{tabular}
\end{center}

The lepton and meson seed triples share a vacuum energy around 
310--350~MeV, squarely in the QCD scale range. This is the 
sBootstrap connection made quantitative: \textbf{the lepton triple 
and the meson seed triple operate at the same scale.}

%======================================================================
\section{A Tale of Two Tuples}
\label{sec:two_tuples}
%======================================================================

\subsection{From seed to full triple}

The seed triples grow into full triples upon symmetry breaking:

\begin{center}
\begin{tabular}{lccc}
\toprule
Sector & Seed triple & Full triple & $v_0^2$ (full) \\
\midrule
Quarks & $(0, s, c)$ & $(-s, c, b)$ & 913 MeV \\
Mesons & $(0, \pi, D_s)$ & $(-\pi, D_s, B)$ & 1230 MeV \\
Leptons & $(0 \approx e, \mu, \tau)$ & $(e, \mu, \tau)$ & 314 MeV \\
\bottomrule
\end{tabular}
\end{center}

In each case:
\begin{itemize}
\item The seed has a massless (or near-massless) member, two massive 
members in the ratio $2 - \sqrt{3}$, and $Q \approx 2/3$.
\item The full triple has \textbf{no} massless member: the zero has 
been lifted (to $m_e$, or the sign-flipped strange/pion mass). The 
$b$ quark and $B$ meson are always present as flavor degrees of 
freedom; what the breaking does is give them their \emph{distinctive 
masses}---the specific values that complete the Koide triple. Before 
breaking, the $b$-direction sits at the flat direction ($\zch = 0$); 
after, it acquires $\zch_b = 64.65\MeV^{1/2}$, the value selected 
by the energy balance with $c$ and the sign-flipped $s$. The bloom 
is not particle creation but \emph{mass generation}.
\end{itemize}

The ``tale of two tuples'' is this: \textbf{each sector presents 
first as a seed with a massless member, then ``blooms'' into a 
full triple when SUSY breaks.} The massless member is the necessary 
initial condition; the breaking lifts it and generates the full spectrum.

\subsection{Precision summary}

\begin{center}
\begin{tabular}{lccc}
\toprule
Triple & $Q$ & Dev.\ from $2/3$ & $v_0^2$ (MeV) \\
\midrule
\multicolumn{4}{c}{\textit{Seed triples}} \\
$(0, \pi, D_s)$ & 0.6679 & 0.18\% & 351 \\
$(0, s, c)$ & 0.6644 & 0.35\% & 228 \\
$(0 \approx e, \mu, \tau)$ & 0.6667 & 0.001\% & 314 \\
\midrule
\multicolumn{4}{c}{\textit{Full triples}} \\
$(e, \mu, \tau)$ & 0.66666 & 0.001\% & 314 \\
$(-\pi, D_s, B)$ & 0.6674 & 0.10\% & 1230 \\
$(-s, c, b)$ & 0.6750 & 1.24\% & 913 \\
$(c, b, t)$ & 0.6695 & 0.42\% & 29\,580 \\
\bottomrule
\end{tabular}
\end{center}

The lepton triple is unique: seed and full triple are nearly 
identical, because the electron mass is small enough to serve as 
the ``zero'' while being nonzero enough to make $Q$ extremely 
precise.

\subsection{Orthogonality of seed and full triples}

Each Koide triple defines a perturbation vector 
$\vec{z} = (z_1, z_2, z_3)$ constrained to the plane $\sum z_k = 0$ 
with length $|\vec{z}|^2 = 3z_0^2$. The angle between seed and 
full triples measures how much the breaking rotates the spectrum 
in charge space:

\begin{center}
\begin{tabular}{lc}
\toprule
Pair of triples & Angle \\
\midrule
$(0,s,c)$ vs.\ $(-s,c,b)$ & $22^\circ$ \\
$(0,\pi,D_s)$ vs.\ $(-\pi,D_s,B)$ & $26^\circ$ \\
$(e,\mu,\tau)$ vs.\ $(0,\mu,\tau)$ & $0.8^\circ$ \\
$(0,s,c)$ vs.\ $(0,\pi,D_s)$ & $0.3^\circ$ \\
\bottomrule
\end{tabular}
\end{center}

Two results are notable. First, the quark seed $(0,s,c)$ and meson 
seed $(0,\pi,D_s)$ are \textbf{nearly parallel} ($0.3^\circ$): the 
meson seed is the composite image of the quark seed, with the GOR 
map preserving direction while changing scale. Second, the lepton 
triple is essentially its own seed ($0.8^\circ$)---the bloom has 
barely begun.


%======================================================================
\section{The Meson Koide Triple: $(-\pi, D_s, B)$}
\label{sec:mesonkoide}
%======================================================================

With the pion charge negative, the triple $(-\pi, D_s, B)$ satisfies:
\begin{equation}
Q(-\pi, D_s, B) = 
\frac{m_\pi + m_{D_s} + m_B}{(-\sqrt{m_\pi} + \sqrt{m_{D_s}} + \sqrt{m_B})^2}
= 0.6674\,,
\end{equation}
deviating from $2/3$ by $0.10\%$---the second most precise Koide 
triple known. The energy balance:
\begin{equation}
\frac{\langle z_k^2\rangle}{z_0^2} = 1.002\,.
\end{equation}

For comparison, the ``standard'' all-positive quark triples $(d,s,b)$ 
and $(u,c,t)$ deviate from $Q = 2/3$ by 9.7\% and 27\% respectively 
at $\overline{\text{MS}}$ masses. \textbf{The meson triple is more 
precise than any quark triple} except $(c,b,t)$.

The pion's negative charge has physical meaning: as a pseudo-Goldstone 
boson, the pion would be massless ($\zch = 0$) if chiral symmetry were 
exact. It sits just below zero ($\zch_\pi = -11.81$), approaching zero 
from the negative side as $m_q \to 0$. The sign of $\zch$ encodes 
Goldstone nature: pseudo-Goldstone bosons sit on the negative side 
of the critical point, having crossed through zero from positive 
territory during chiral symmetry breaking.

%======================================================================
\section{The Third Tuple: $(c, b, t)$}
\label{sec:cbt}
%======================================================================

\subsection{A tuple without mesons}

The triple $(c, b, t)$ with all positive signs gives $Q = 0.6695$, 
deviating by $0.42\%$. Its vacuum energy $v_0^2 \approx 29.6\GeV$ 
is far above the QCD scale.

This triple has \textbf{no meson counterpart}: the top quark decays 
weakly ($t \to Wb$, lifetime $\sim 5 \times 10^{-25}$~s) before the 
strong interaction can bind it ($\tau_{\text{QCD}} \sim 1/\Lambda_{\text{QCD}} 
\sim 10^{-24}$~s). No $T$-meson ($t\bar{q}$) exists. The meson matrix 
$M^i{}_j$ extends only to $b$-flavored states; the top sector lives 
purely at the quark level.

\subsection{Diquarks and the top}

In the sBootstrap, diquarks $[q_i q_j]$ are the scalar partners of 
antiquarks $\bar{q}_k$, with the baryon acting as the SUSY generator 
\cite{CattoGursey, Miyazawa}. For five active flavors 
$(u, d, s, c, b)$, the diquarks form a representation of $SU(5)_f$.

The top diquarks $[tq]$ are kinematic objects only---they do not 
hadronize. Yet their masses should follow the same pattern as other 
diquarks: \textbf{there is no dynamical reason for top diquark masses 
to differ from the others}. The $(c,b,t)$ Koide triple governs the 
quark masses that enter the diquark sector, and its energy balance 
$\langle z_k^2\rangle = z_0^2$ constrains the possible diquark 
spectrum even though the top diquarks are never observed as bound 
states.

\subsection{The $(c,b,t)$ triple has no seed}

Unlike the $(0,s,c)$ and $(0,\pi,D_s)$ seeds, there is no obvious 
$(0, x, y)$ seed for the $(c,b,t)$ triple. No pair of quarks in this 
sector satisfies $\sqrt{m_{\text{light}}}/\sqrt{m_{\text{heavy}}} 
\approx 2 - \sqrt{3}$: one has $\sqrt{m_c}/\sqrt{m_t} = 0.086$ 
(too small) and $\sqrt{m_c}/\sqrt{m_b} = 0.551$ (too large).

This means $(c,b,t)$ is \emph{not} a ``bloomed seed'' in the same 
sense as $(-s,c,b)$. It may represent a fundamentally different 
structure---a triple that was never seeded by a massless member but 
instead emerged fully formed at a higher breaking scale, perhaps 
connected to electroweak symmetry breaking rather than to chiral 
symmetry breaking.

%======================================================================
\section{The Diquark Sector and SU(5) Flavor}
\label{sec:diquarks}
%======================================================================

\subsection{Counting partners}

In the full sBootstrap \cite{Rivero2005}, the scalar partners of 
the Standard Model fermions include both mesons ($q\bar{q}$) and 
diquarks ($qq$). For five active flavors, the meson matrix gives 
$5 \times 5 = 25$ scalars. The diquarks $[q_i q_j]$ form a 
representation of $SU(5)_f$:
\begin{equation}
\mathbf{5} \otimes \mathbf{5} = \overline{\mathbf{10}}_A 
\oplus \mathbf{15}_S\,.
\end{equation}

The symmetric $\mathbf{15}_S$ gives 15 diquarks including 
same-flavor states $[uu]$, $[dd]$, \ldots, while the antisymmetric 
$\overline{\mathbf{10}}_A$ gives 10 diquarks with $i \neq j$ only.

\subsection{Why antisymmetric?}

For the sBootstrap to produce the exact count of scalar partners 
matching the Standard Model fermion content, the diquarks must 
transform in the \textbf{antisymmetric} $\overline{\mathbf{10}}$ 
\cite{Rivero2005}. The argument:

The 25 mesons account for the slepton and ``internal'' squark 
sectors. The remaining squarks must come from diquarks. The 
$\overline{\mathbf{10}}$ provides exactly 10 diquarks: $[ud]$, 
$[us]$, $[uc]$, $[ub]$, $[ds]$, $[dc]$, $[db]$, $[sc]$, $[sb]$, 
$[cb]$. With their antiparticles, this gives 20 scalar states. The 
total $25 + 10 = 35$ (or $25 + 20 = 45$ counting antiparticles 
separately) matches the scalar content of the Supersymmetric 
Standard Model.

The symmetric $\mathbf{15}$ would give 15 diquarks, overcounting 
by 5 (the same-flavor states). Moreover, the exotic $+4/3$ charged 
diquarks $[uu]$, $[uc]$, $[cc]$ have no lepton partner among known 
particles---a fatal obstacle.

\subsection{The Pauli problem}

The antisymmetric flavor requirement conflicts with a naive 
application of Fermi statistics. For spin-0, $L=0$ diquarks in 
the color $\bar{3}$ (antisymmetric) channel:

\begin{center}
\begin{tabular}{lcl}
\toprule
Property & Symmetry & Wavefunction \\
\midrule
Color $\bar{3}$ & Antisymmetric & $\epsilon_{abc} q^b q^c$ \\
Spin 0 & Antisymmetric & $(\uparrow\downarrow - \downarrow\uparrow)/\sqrt{2}$ \\
Spatial $L=0$ & Symmetric & $\psi(r)$ \\
\midrule
Overall & Must be A & $A \times A \times S = S$ \quad$(\neq A!)$ \\
\bottomrule
\end{tabular}
\end{center}

The product is symmetric---\textbf{flavor must also be symmetric} 
to make the overall wavefunction antisymmetric. But the sBootstrap 
requires antisymmetric flavor.

This is one of the deepest puzzles of the program. Possible 
resolutions:
\begin{itemize}
\item \textbf{$L = 1$ diquarks:} Adding orbital angular momentum 
introduces spatial antisymmetry. Then color(A) $\times$ spin(A) 
$\times$ space(A) $\times$ flavor = overall A requires 
flavor(S)---still wrong. But $L=1$ with spin-1 gives spin(S), 
and then color(A) $\times$ spin(S) $\times$ space(A) $\times$ 
flavor(?) = A requires flavor(A). This works but requires both 
$L=1$ and spin-1 diquarks, departing from the simplest scalar 
diquark picture.
\item \textbf{Dynamical projection:} A binding mechanism that 
effectively projects onto the flavor-antisymmetric channel, perhaps 
through instanton effects or a condensate in the breaking sector.
\item \textbf{The SO(32) route:} In the string-theoretic embedding 
\cite{Rivero2005}, the $SO(32)$ spinor representation automatically 
produces the correct antisymmetric structure. The Pauli constraint 
is an artifact of treating diquarks as simple two-quark states; in 
the string picture, they are fundamental objects of the open-string 
sector and the symmetry is determined by the gauge group, not by 
quark statistics.
\end{itemize}

\subsection{The SO(32) connection}

The sBootstrap counting matches a decomposition of the $SO(32)$ gauge 
group familiar from Type~I/heterotic string theory. Under 
$SO(32) \supset SU(5) \times \ldots$, the adjoint or spinor 
representations produce exactly the right combination of mesons 
(in $\mathbf{5} \otimes \bar{\mathbf{5}}$) and antisymmetric 
diquarks (in $\overline{\mathbf{10}}$). The 496-dimensional adjoint 
of $SO(32)$ decomposes to include the 24 of $SU(5)$ (gauge bosons), 
the $\mathbf{5} \otimes \bar{\mathbf{5}} = \mathbf{1} + \mathbf{24}$ 
(mesons/sleptons), and the $\overline{\mathbf{10}} + \mathbf{10}$ 
(diquarks/squarks).

This is suggestive: the same structure that gives anomaly-free string 
compactifications also gives the correct sBootstrap particle count 
with the required antisymmetric diquarks. Whether this is a 
coincidence or a hint of UV completion is an open question.

%======================================================================
\section{SUSY Breaking: From Seeds to Spectrum}
\label{sec:breaking}
%======================================================================

The SUSY-breaking story begins with the foundational observation 
\eqref{eq:pimu}: $m_\pi^2 - m_\mu^2 \approx f_\pi^2$. This is the 
Volkov--Akulov mass formula for a pseudo-Goldstino---the muon---whose 
superpartner---the pion---is a pseudo-Goldstone boson. The SUSY-breaking 
scale is $f_\pi \approx 92\MeV$, and the breaking mechanism is chiral: 
$\qqbar \neq 0$ breaks both chiral symmetry and the sBootstrap SUSY 
simultaneously. The Volkov--Akulov realization is the \emph{physical 
starting point}; what follows in this section is the SQCD machinery 
that makes it precise.

The dynamics must accomplish three things: generate the seed-to-bloom 
transition for each tuple, preserve the energy-balance condition 
$\langle z_k^2\rangle = z_0^2$ through the breaking, and produce 
the correct splitting between meson and lepton masses at scale $f_\pi$.

\subsection{The Seiberg framework}

The natural arena is $\mathcal{N}=1$ SQCD with gauge group $SU(3)_c$ 
and $N_f$ flavors of chiral superfields $Q^i$, $\bar{Q}_j$ in the 
fundamental and antifundamental representations. For $N_f = N_c = 3$ 
(the case relevant to light QCD), Seiberg's exact results 
\cite{Seiberg1995} establish that the theory confines, with the 
low-energy degrees of freedom being the composite meson superfield
\begin{equation}
M^i{}_j = Q^i \cdot \bar{Q}_j\,,
\end{equation}
a $3\times 3$ matrix of chiral superfields, together with baryons 
$B = \epsilon_{ijk}Q^iQ^jQ^k$ and $\bar{B}$. These satisfy the 
quantum-modified constraint
\begin{equation}
\label{eq:seiberg}
\det M - B\bar{B} = \Lambda^{2N_c}\,,
\end{equation}
where $\Lambda$ is the dynamical scale.

Each entry $M^i{}_j$ is a chiral superfield containing a complex 
scalar (the meson) and a Weyl fermion (the lepton, in the sBootstrap 
identification). The scalar and fermion are degenerate in the SUSY 
limit; SUSY breaking splits them.

\subsection{The three chiral superfields of each seed}

The seed triple $(0, \zch_a, \zch_b)$ maps naturally onto three 
chiral superfields with the following mass structure before breaking:

\begin{center}
\begin{tabular}{lcl}
\toprule
Superfield & Mass & Role \\
\midrule
$\Phi_0$ & $m_0 = 0$ & Massless: flat direction \\
$\Phi_a$ & $m_a = \zch_a^2$ & Light massive member \\
$\Phi_b$ & $m_b = \zch_b^2$ & Heavy massive member \\
\bottomrule
\end{tabular}
\end{center}

The ratio $\zch_a/\zch_b = 2 - \sqrt{3}$ is enforced by the 
Koide energy-balance condition applied to the seed.

In the Seiberg description, these three superfields are entries of 
the meson matrix $M^i{}_j$. For the $(0, s, c)$ seed, the relevant 
entries involve the strange and charm quarks; $\Phi_0$ is the entry 
whose quark-mass combination vanishes in the seed limit (i.e., the 
direction in flavor space where $m_i + m_j \to 0$ through GOR).

\subsection{F-term breaking: the O'Raifeartaigh mechanism}
\label{subsec:fterm}

The VA relation \eqref{eq:pimu} tells us \emph{that} SUSY is broken at 
scale $f_\pi$. To understand \emph{how}, we need the microscopic 
mechanism. The O'Raifeartaigh model \cite{ORaifeartaigh} is the simplest 
realization of spontaneous $F$-term SUSY breaking. With three chiral 
superfields, the superpotential
\begin{equation}
W = f\,\Phi_0 + m\,\Phi_1\Phi_2 + g\,\Phi_0\Phi_1^2
\end{equation}
gives $F$-terms:
\begin{align}
F_{\Phi_0}^* &= f + g\,\phi_1^2\,, \\
F_{\Phi_1}^* &= m\,\phi_2 + 2g\,\phi_0\phi_1\,, \\
F_{\Phi_2}^* &= m\,\phi_1\,.
\end{align}
If $f \neq 0$ and $m \neq 0$, then $F_{\Phi_2}^* = 0$ requires 
$\phi_1 = 0$, but then $F_{\Phi_0}^* = f \neq 0$: SUSY is 
\textbf{necessarily broken}. The key features are:
\begin{enumerate}
\item A massless field ($\Phi_0$) exists at tree level---it is the 
pseudo-modulus.
\item The breaking is triggered because $F_{\Phi_0}$ cannot vanish.
\item After including loop corrections, $\Phi_0$ acquires a mass 
(the Coleman--Weinberg potential lifts the flat direction).
\end{enumerate}

The structural parallel with the seed triples is:

\begin{center}
\begin{tabular}{llll}
\toprule
O'Raifeartaigh & Flavor & Seed & After breaking \\
\midrule
$\Phi_0$ (massless) & $b$ / $B$ & $\zch = 0$ (flat direction) & $\zch_b = 64.65$ (mass generated) \\
$\Phi_1$ (light) & $s$ / $\pi$ & $\zch_s = 9.67$ & $\zch \to -9.67$ (sign-flipped) \\
$\Phi_2$ (heavy) & $c$ / $D_s$ & $\zch_c = 35.64$ & $\zch_c \approx 35.64$ (shifted) \\
\midrule
 & & $f = 0$ (SUSY) & $f = f_\pi$ (broken) \\
\bottomrule
\end{tabular}
\end{center}

This parallel is \emph{structural}, not quantitative: the minimal 
O'Raifeartaigh model gives $\Phi_1$ and $\Phi_2$ the same mass 
parameter $m$, whereas the seed triple has a mass ratio 
$m_s/m_c = (2 - \sqrt{3})^2 \approx 0.07$---far from degenerate. 
The hierarchy between the two massive members must come from 
additional dynamics beyond the minimal O'Raifeartaigh superpotential. 
What the O'Raifeartaigh structure \emph{does} provide is the essential 
feature: a massless flat direction ($\Phi_0 = $ the $b$-direction) 
whose $F$-term is forced nonzero. The $b$ quark and $B$ meson are 
always present as flavor degrees of freedom; the seed configuration 
$(0, s, c)$ means the $b$-direction sits at the flat direction before 
breaking. The bloom is not particle creation but mass generation: 
breaking lifts $\Phi_0$ from zero to $\zch_b \neq 0$, simultaneously 
flipping the sign of $\zch_s$, so that the same three flavors 
rearrange from $(0, s, c)$ into $(-s, c, b)$.

The identification $f = \langle F_{\Phi_0}\rangle \sim f_\pi$ is the 
sBootstrap's central claim: the SUSY-breaking $F$-term is the chiral 
condensate parameter.

In SQCD, this is realized dynamically rather than through a 
superpotential: the chiral condensate $\qqbar \neq 0$ acts as the 
nonzero $F$-term, the impossibility of restoring chiral symmetry 
in the QCD vacuum is the dynamical analogue of $f \neq 0$, and the 
pion---the would-be Goldstone boson that acquires mass from explicit 
chiral breaking---plays the role of $\Phi_0$ whose mass is lifted 
by loop corrections (here, by quark masses).

\subsection{The Goldstino direction and the Volkov--Akulov realization}
\label{subsec:va}

When SUSY is spontaneously broken, a massless fermion appears: the 
Goldstino, the fermionic analogue of the Goldstone boson. Its decay 
constant $f_{\text{VA}}$ sets the SUSY-breaking scale through the 
Volkov--Akulov Lagrangian \cite{VolkovAkulov1973}:
\begin{equation}
\mathcal{L}_{\text{VA}} = -f_{\text{VA}}^2 \det\!\left(
\delta_\mu^\nu + \frac{i}{f_{\text{VA}}^2}\,\bar{\lambda}\bar{\sigma}^\nu
\partial_\mu\lambda\right),
\end{equation}
where $\lambda$ is the Goldstino field. At low energies, the Goldstino 
couples derivatively with strength $\sim 1/f_{\text{VA}}^2$, exactly 
as the pion couples with strength $\sim 1/f_\pi^2$ in the chiral 
Lagrangian.

The sBootstrap identification is:
\begin{equation}
\label{eq:va_id}
f_{\text{VA}} = f_\pi\,, \qquad 
\text{Goldstino} = \text{muon (partner of }\pi\text{)}\,.
\end{equation}
More precisely, the muon is a \emph{pseudo-Goldstino}: in the 
limit $m_q \to 0$ where the pion becomes exactly massless (a true 
Goldstone boson), its fermionic partner also becomes massless (a 
true Goldstino). The finite quark masses explicitly break both 
chiral symmetry (giving $m_\pi \neq 0$) and the sBootstrap SUSY 
(giving $m_\mu \neq m_\pi$), with the splitting controlled by 
$f_\pi$:
\begin{equation}
m_\pi^2 - m_\mu^2 = f_\pi^2\,.
\end{equation}

This is the Volkov--Akulov relation for the pseudo-Goldstino: the 
mass splitting between the scalar (pion) and its fermionic partner 
(muon) is set by the SUSY-breaking scale squared.

The VA mechanism is not a separate breaking from the F-term: it is 
the \emph{low-energy, composite-level description} of the same 
breaking. At the microscopic level, $\langle F\rangle \neq 0$ 
in the meson superfield; at the composite level, this manifests 
as the Goldstino coupling with $f_{\text{VA}} = f_\pi$. These 
two quantities---the Goldstino decay constant and the pion decay 
constant---are \emph{identified} in the sBootstrap; this is the 
central claim, not a derivation.

The identification should not be confused with the standard relation 
between the $F$-term VEV and the chiral condensate. In standard 
SQCD, $\langle F\rangle \sim |\qqbar|/f_\pi \sim (250\MeV)^3/92\MeV 
\approx 170{,}000\MeV^2$, while $f_\pi^2 \approx 8{,}500\MeV^2$---a 
factor of $\sim 20$ discrepancy. The GOR relation 
$m_\pi^2 f_\pi^2 = -(m_u + m_d)\qqbar$ connects these through the 
small quark masses, but does not equate them. The sBootstrap 
identification $f_{\text{VA}} = f_\pi$ is therefore a 
\emph{low-energy effective} statement about the composite sector, 
not a microscopic identity. It says: the scale at which SUSY 
appears broken in the meson--lepton spectrum is $f_\pi$, regardless 
of the microscopic $F$-term that produces this breaking.

The VA realization gives more than a mass formula---it organizes the 
entire breaking narrative. In the chiral limit ($m_q \to 0$), both 
pion and muon are exactly massless: the pion as the Goldstone boson 
of broken chiral symmetry, the muon as the Goldstino of broken 
sBootstrap SUSY. This is the \emph{exact seed}: the $B$-meson 
direction sits at the flat direction ($\zch = 0$), and the 
$\pi$--$\mu$ pair is massless. Turning on quark masses ($m_q \neq 0$) 
does three things simultaneously:
\begin{enumerate}
\item Gives the pion its GOR mass (lifting the flat direction for $\pi$).
\item Gives the muon its mass (the pseudo-Goldstino acquires mass 
when the explicit breaking parameter $m_q$ is nonzero).
\item Generates the $B$-meson mass (the flat direction for $\Phi_0$ 
is lifted by the Coleman--Weinberg potential, i.e., by the same 
quark masses that feed the GOR formula).
\end{enumerate}
The splitting between (1) and (2) is $f_\pi^2$---the VA relation. 
The seed-to-bloom transition of \S\ref{sec:two_tuples} is the same 
physics seen from the charge-space perspective: $m_q \neq 0$ lifts 
the zero, generates the $b$/$B$ mass, flips the sign of $\zch_s$ 
and $\zch_\pi$, and produces the full triple. The VA equation 
\eqref{eq:pimu} is not one fact among many; it is the generating 
relation from which the bloom follows.

A note on quantum numbers: the Goldstino of spontaneously broken 
spacetime SUSY is a neutral fermion (the fermionic component of the 
superfield whose $F$-term breaks SUSY). The muon carries electric 
charge and lepton number. The sBootstrap SUSY, however, is not a 
spacetime symmetry but an emergent hadronic symmetry---a 
Miyazawa-type superalgebra \cite{Miyazawa} acting on composite 
states. In this context, the ``Goldstino'' need not carry the 
quantum numbers of a spacetime supercharge; it carries the quantum 
numbers appropriate to the broken hadronic generator. The muon's 
charge is inherited from the pion multiplet it partners.

\subsection{D-term breaking: electromagnetic splittings}
\label{subsec:dterm}

The unbroken $U(1)_{\text{em}}$ gauge symmetry allows a 
Fayet--Iliopoulos (FI) term \cite{FayetIliopoulos}:
\begin{equation}
\mathcal{L}_{\text{FI}} = \xi\,D_{\text{em}}\,,
\end{equation}
which shifts scalar masses by $\delta m^2 = Q_{\text{em}}\,\xi$, 
where $Q_{\text{em}}$ is the electric charge. This is the 
supersymmetric version of the electromagnetic mass difference.

In QCD without SUSY, the electromagnetic mass splitting of the 
pion is given by the Das--Guralnik--Mathur--Low--Young (DGMLY) 
formula \cite{DGMLY1967}:
\begin{equation}
m_{\pi^\pm}^2 - m_{\pi^0}^2 = \frac{3\alpha_{\text{em}}}{4\pi}\,
\frac{m_\rho^2}{f_\pi^2}\,f_\pi^2 
\approx (35.5\MeV)^2\,,
\end{equation}
which in the SUSY framework becomes a D-term contribution:
\begin{equation}
\xi \sim \frac{3\alpha_{\text{em}}}{4\pi}\,m_\rho^2 
\sim (35\MeV)^2\,.
\end{equation}

In charge space, the D-term shifts $\zch$ by an amount proportional 
to $Q_{\text{em}}$:
\begin{equation}
\zch \to \zch + \delta\zch\,, \qquad 
\delta\zch \approx \frac{Q_{\text{em}}\,\xi}{2\zch}\,.
\end{equation}
For the pion ($Q_{\text{em}} = \pm 1$, $\zch \approx 11.8$): 
$\delta\zch \approx 0.2\MeV^{1/2}$, matching the observed 
$\zch_{\pi^\pm} - \zch_{\pi^0} = 11.82 - 11.62 = 0.20\MeV^{1/2}$.

The D-term effect is subleading to the F-term: 
$\xi \sim (35\MeV)^2$ versus $f_\pi^2 \sim (92\MeV)^2$. It 
operates \emph{within} each supermultiplet entry, splitting charged 
from neutral mesons, while the F-term splits bosons from fermions 
across the supermultiplet.

Witten's inequality \cite{Witten1983} guarantees 
$m_{\pi^\pm} > m_{\pi^0}$ in QCD with electromagnetism---the 
charged pion is always heavier. In the SUSY framework this 
becomes the statement that the D-term contribution is positive 
for $Q_{\text{em}} \neq 0$, consistent with the FI sign.

\subsection{Mapping the breaking to the tuples}
\label{subsec:mapping}

The tuple structure maps onto a layered breaking with two independent 
mechanisms (F-term and D-term) plus a separate heavy-sector origin. 
The Volkov--Akulov realization of \S\ref{subsec:va} is not a third 
mechanism: it is the low-energy, composite-level description of the 
same F-term breaking. The VA relation \eqref{eq:pimu} \emph{sets the 
scale} ($f_{\text{VA}} = f_\pi$); the O'Raifeartaigh/SQCD dynamics 
of \S\ref{subsec:fterm} \emph{provides the mechanism} 
($\langle F\rangle \neq 0$).

\paragraph{Layer 1: F-term/VA (chiral condensate, scale $\sim f_\pi$).}
The chiral condensate $\qqbar \neq 0$ breaks both chiral symmetry 
and the sBootstrap SUSY. This is the primary breaking, affecting 
the light sector. The seed triples $(0, s, c)$ and $(0, \pi, D_s)$ 
are the pre-breaking configurations; the massless member is the flat 
direction whose $F$-term is forced nonzero by confinement. After 
breaking:
\begin{itemize}
\item The zero member is lifted and acquires negative charge 
($\zch < 0$), becoming the sign-flipped member of the full triple 
($-s$, $-\pi$).
\item The full triples $(-s,c,b)$ and $(-\pi, D_s, B)$ emerge, 
incorporating the $b$-quark and $B$-meson sectors.
\item The lepton triple $(e, \mu, \tau)$ is the fermionic shadow: 
the leptons acquire masses through the same condensate, with the 
muon as pseudo-Goldstino.
\end{itemize}

\paragraph{Layer 2: D-term (electromagnetism, scale $\sim \alpha_{\text{em}} m_\rho$).}
The $U(1)_{\text{em}}$ FI term generates electromagnetic mass 
splittings within each multiplet. This is a perturbative correction 
on top of the F-term breaking, affecting only charged particles. 
It does not create new Koide triples but shifts existing ones by 
$O(\alpha_{\text{em}})$.

\paragraph{Layer 3: The heavy sector (scale $\sim v_0^{(cbt)}$).}
The $(c,b,t)$ triple at $v_0^2 \approx 29.6\GeV$ operates far 
above the chiral scale. It has no seed (no massless member), no 
meson shadow (top does not hadronize), and its breaking may be 
connected to a different sector of the theory---possibly the 
electroweak Higgs mechanism.

This layered structure involves three distinct scales, each with 
different physical meaning:
\begin{align}
\text{F-term (boson--fermion splitting):} &\quad 
f_\pi \approx 92\MeV\,, \quad f_\pi^2 \approx 8{,}500\MeV^2\,; \\
\text{D-term (EM fine structure):} &\quad 
\xi^{1/2} \approx 35\MeV\,, \quad \xi \approx 1{,}200\MeV^2\,; \\
\text{Heavy sector (Koide scale):} &\quad 
v_0^{(cbt)} \approx 172\MeV^{1/2}\,, \quad (v_0^{(cbt)})^2 \approx 29{,}600\MeV\,.
\end{align}
Comparing the squared scales: $f_\pi^2 \gg \xi$ (F-term dominates 
over D-term by a factor $\sim 7$, consistent with the D-term being 
a perturbative $O(\alpha_{\text{em}})$ correction). The heavy-sector 
scale $(v_0^{(cbt)})^2$ carries dimensions of $\MeV$ (not $\MeV^2$) 
because $v_0$ is a charge ($\MeV^{1/2}$), not an energy; direct 
comparison with $f_\pi$ requires specifying which physical quantity 
is being compared.

\subsection{The O'Raifeartaigh structure in SQCD}

In Seiberg's SQCD with $N_f = N_c = 3$, the quantum constraint 
\eqref{eq:seiberg} provides a natural setting for the 
O'Raifeartaigh mechanism. The constraint $\det M - B\bar{B} = \Lambda^6$ 
means the moduli space has a specific geometry: not all field directions 
are flat. When quark masses are added via a tree-level superpotential 
$W_{\text{tree}} = m_i M^i{}_i$ (sum over flavors), the $F$-terms become:
\begin{equation}
F_{M^i{}_j} = m_i\,\delta^i_j + \Lambda^6\,(\text{constraint terms})\,.
\end{equation}

For the seed triple $(0, s, c)$, the ``$0$'' member corresponds to a 
direction in the meson matrix where the tree-level mass vanishes---the 
entry involving the lightest quark combination. In the exact chiral 
limit ($m_u = m_d = 0$), this entry is truly massless, and the 
constraint \eqref{eq:seiberg} forces $\langle F\rangle \neq 0$ for 
that entry. This is the dynamical O'Raifeartaigh mechanism: 
confinement (encoded in $\Lambda^6$) plays the role of the parameter 
$f$ in the superpotential, and the impossibility of setting all 
$F$-terms to zero is a consequence of the strong dynamics rather 
than a choice of superpotential.

The ISS mechanism (Intriligator--Seiberg--Shih \cite{ISS2006}) 
establishes metastable SUSY-breaking vacua for $N_c < N_f < 3N_c/2$, 
where the rank condition in the magnetic dual prevents all $F$-terms 
from vanishing. The sBootstrap case $N_f = N_c = 3$ is strictly 
\emph{below} the ISS window---at $N_f = N_c$ the low-energy 
description is qualitatively different, governed by the 
quantum-modified constraint \eqref{eq:seiberg} on the moduli space 
rather than by a free magnetic dual. However, the physical QCD has 
$N_f = 3$ light flavors and $N_f = 4$--$5$ if charm and bottom are 
included at their respective thresholds, which does fall within the 
ISS window for $N_c = 3$. The sBootstrap's multi-scale structure---with 
the seed triples involving $s$, $c$, and $b$ quarks at different 
scales---may require treating the effective $N_f$ as 
scale-dependent, entering the ISS regime above the strange threshold.

\subsection{How the three supermultiplets join}

The seed-to-bloom transition generates three separate Koide triples. 
But the full sBootstrap requires these triples to be \emph{embedded 
in a single meson matrix}. The $3\times 3$ light-flavor matrix is:

\begin{equation}
M^{(3)} = \begin{pmatrix}
u\bar{u} & u\bar{d} & u\bar{s} \\
d\bar{u} & d\bar{d} & d\bar{s} \\
s\bar{u} & s\bar{d} & s\bar{s}
\end{pmatrix}
\;\longleftrightarrow\;
\begin{pmatrix}
\pi^0/\eta & \pi^+ & K^+ \\
\pi^- & \pi^0/\eta & K^0 \\
K^- & \bar{K}^0 & \eta'
\end{pmatrix}.
\end{equation}

When charm and bottom are included, the matrix extends to $5\times 5$. 
We display it schematically, with the members of the meson Koide 
triple $(-\pi, D_s, B)$ marked by boxes:

\begin{equation}
M^{(5)} = \begin{pmatrix}
\pi^0/\eta & \boxed{\pi^+} & K^+ & D^0 & \boxed{B^+} \\
\pi^- & \pi^0/\eta & K^0 & D^- & B^0 \\
K^- & \bar{K}^0 & \eta' & \boxed{D_s^-} & B_s \\
\bar{D}^0 & D^+ & D_s^+ & \eta_c & B_c \\
B^- & \bar{B}^0 & \bar{B}_s & \bar{B}_c & \eta_b
\end{pmatrix},
\end{equation}
where rows are indexed by $(u, d, s, c, b)$ and columns by 
$(\bar{u}, \bar{d}, \bar{s}, \bar{c}, \bar{b})$.

The three members of the meson Koide triple occupy different 
off-diagonal blocks of this matrix: $\pi^+ = u\bar{d}$ is in the 
light $(u,d)$ block, $D_s^- = s\bar{c}$ bridges the light-charm 
sector, and $B^+ = u\bar{b}$ bridges the light-bottom sector. They 
do not form a ``slice'' in any simple linear-algebraic sense---they 
are scattered entries connected by the Koide condition on their 
\emph{masses}, not by their positions in the matrix. The meson seed 
$(0, \pi, D_s)$ drops the $B^+$ and replaces it with zero (the 
entry that would involve the top quark, absent from the matrix).

This scattered structure is the central puzzle: the Koide triples 
are \emph{mass relations} among matrix entries, not submatrices. 
The question is whether the Seiberg constraint 
$\det M - B\bar{B} = \Lambda^6$ (appropriately generalized to 
$5 \times 5$), combined with the SUSY-breaking dynamics, produces 
correlations among these scattered entries that enforce the 
energy-balance condition $\langle z_k^2\rangle = z_0^2$.

\subsection{The energy balance as a fixed point}

One possibility is that the energy-balance condition 
$\langle z_k^2\rangle = z_0^2$ is a \emph{dynamical attractor}: 
any set of three masses evolving under the SUSY-breaking dynamics 
flows toward the Koide surface $Q = 2/3$, as a fixed point of the 
RG flow for the soft SUSY-breaking parameters. This is speculative, 
but the observation that the condition holds for fundamental fermions 
(leptons, quarks) \emph{and} composite bosons (mesons) across very 
different scales lends it plausibility. A concrete test would be to 
show that the one-loop soft-mass RG equations in $\mathcal{N}=1$ 
SQCD with appropriate boundary conditions flow toward 
$\sigma(\zch)/\bar{\zch} = 1$; this calculation has not yet been 
performed.

\subsection{The exact seed and what breaks it}
\label{subsec:exactseed}

The seed configuration $(0, \zch_a, \zch_b)$ requires \emph{two} 
independent conditions:
\begin{enumerate}
\item One member at $\zch = 0$ exactly (the flat direction).
\item The ratio of the massive members equals $2 - \sqrt{3}$:
\begin{equation}
\label{eq:seedratio}
\frac{\sqrt{m_a}}{\sqrt{m_b}} = 2 - \sqrt{3} \approx 0.26795\,.
\end{equation}
\end{enumerate}
These are logically distinct. The chiral limit ($m_u = m_d = 0$, 
$\qqbar \neq 0$) enforces condition (1) for the meson sector: the 
pion is an exact Goldstone boson, $m_\pi = 0$, and in the 
picture of \S\ref{subsec:fterm} the $B$-meson direction sits at 
the flat direction before acquiring mass. But the chiral limit says 
\emph{nothing} about condition (2)---it does not constrain 
$\sqrt{m_{D_s}}/\sqrt{m_B}$ or $\sqrt{m_s}/\sqrt{m_c}$ 
to equal $2 - \sqrt{3}$.

The physical deviations from the exact ratio are small but 
measurable:

\begin{center}
\begin{tabular}{lccc}
\toprule
Seed & Exact ratio & Physical ratio & Deviation \\
\midrule
$(0, \pi, D_s)$ & 0.2680 & 
$\sqrt{m_\pi}/\sqrt{m_{D_s}} = 0.2663$ & $-0.6\%$ \\
$(0, s, c)$ & 0.2680 & 
$\sqrt{m_s}/\sqrt{m_c} = 0.2712$ & $+1.2\%$ \\
$(0 \approx e, \mu, \tau)$ & 0.2680 & 
$\sqrt{m_\mu}/\sqrt{m_\tau} = 0.2439$ & $-9.0\%$ \\
\bottomrule
\end{tabular}
\end{center}

The meson seed is closest to the exact ratio (0.6\% deviation). 
This is expected: the meson masses are the most ``infrared'' 
quantities, closest to the chiral dynamics that creates the seed. 
The quark seed deviates more (1.2\%) because the quark masses 
$m_s$ and $m_c$ are $\overline{\text{MS}}$ running masses, 
evaluated at $\mu = 2\GeV$. As the renormalization scale $\mu$ 
varies, the ratio $\sqrt{m_s(\mu)}/\sqrt{m_c(\mu)}$ changes 
because $m_s$ and $m_c$ run at different rates (different anomalous 
dimensions in QCD). This raises a natural question:

\emph{Does there exist a scale $\mu^*$ at which the ratio 
$\sqrt{m_s(\mu^*)}/\sqrt{m_c(\mu^*)}$ equals $2 - \sqrt{3}$ 
exactly?}

If so, $\mu^*$ is the ``seed scale''---the energy at which the 
exact seed configuration is realized. Since $m_s$ runs faster 
than $m_c$ (the anomalous dimension is larger for the lighter 
quark), the ratio $\sqrt{m_s}/\sqrt{m_c}$ decreases with 
increasing $\mu$. At $\mu = 2\GeV$ the ratio is 0.2712 (above 
the target 0.2680); it will reach 0.2680 at some 
$\mu^* > 2\GeV$, likely in the range 3--5 GeV. This is a 
concrete, testable prediction.

The lepton ``seed'' has the worst ratio match ($-9\%$), which 
is consistent with the observation that the lepton triple is 
\emph{not} a true seed: $m_e \neq 0$, and the electron was never 
at the flat direction. The leptons are an ``almost-seed'' whose 
excellent $Q = 0.66666$ arises not from the seed structure but 
from the smallness of $m_e$, which makes the two conditions 
approximately degenerate. For the leptons, condition (1) is 
approximate rather than exact, and condition (2) is correspondingly 
relaxed.

The physical origin of the deviations from condition (2) is 
QCD dynamics: running of quark masses, higher-order chiral 
perturbation theory corrections to meson masses, and 
nonperturbative effects. In principle, these are calculable. 
A next-to-leading-order chiral perturbation theory calculation 
of $m_{D_s}$ with the Koide constraint could determine 
whether the 0.6\% meson deviation is consistent with known 
chiral corrections, or whether it requires new physics.

\subsection{What the $(c,b,t)$ triple tells us}

The $(c,b,t)$ triple ($Q = 0.6695$, $v_0^2 = 29.6\GeV$) has 
neither a seed nor a meson shadow. It operates at a scale far 
above $\Lambda_{\text{QCD}}$, suggesting its origin is different:

\begin{enumerate}
\item It may be governed by the electroweak Higgs mechanism rather 
than the chiral condensate. The top Yukawa $y_t \approx 1$ is 
``natural'' in the sense that $m_t \approx v_{\text{EW}}/\sqrt{2}$; 
the Koide condition on $(c,b,t)$ would then constrain the Yukawa 
coupling matrix.
\item It may represent a Volkov--Akulov realization at the 
electroweak scale: if there is a second SUSY breaking at 
$f' \sim v_{\text{EW}}$, the $(c,b,t)$ triple would be its 
analog of the $(\pi, \mu)$ system, but with $f'$ replacing $f_\pi$.
\item Or it may be a purely kinematic accident---the quark masses 
happen to satisfy the Koide condition without a dynamical reason. 
The 0.42\% precision makes this unlikely but not impossible.
\end{enumerate}

Interpretation (1) is most natural within the sBootstrap: the 
$(c,b,t)$ triple involves quarks whose masses are dominated by 
Yukawa couplings to the Higgs, not by the chiral condensate, and the 
Koide condition would then be a constraint on the Yukawa matrix 
rather than on the SQCD dynamics. The chiral sector and the 
electroweak sector would each independently produce a Koide-type 
energy balance, with the charm and bottom quarks serving as the 
bridge between the two.

If interpretation (2) is also invoked, the sBootstrap has two breaking 
scales:
\begin{equation}
f_{\text{chiral}} = f_\pi \approx 92\MeV\,, \qquad
f_{\text{EW}} \sim v_{\text{EW}} \approx 246\GeV\,,
\end{equation}
and the full spectrum requires both. The light sector (leptons, 
light mesons) is governed by $f_\pi$; the heavy sector (top, and 
by extension charm and bottom through the Koide network) is 
governed by $v_{\text{EW}}$. The charm and bottom quarks, 
appearing in both the $(0,s,c)/(-s,c,b)$ light-sector triples 
and the $(c,b,t)$ heavy-sector triple, are the ``bridges'' 
between the two scales.

%======================================================================
\section{The Electron Mass and the Lepton Seed}
\label{sec:electron}
%======================================================================

The charged lepton triple $(e, \mu, \tau)$ is the ``almost-seed'': 
$\zch_e = 0.715\MeV^{1/2}$ is small but nonzero. The electron's 
charge is near zero because of cancellation in $\zch_e = z_0 + z_e$:
\begin{equation}
v_0 = 17.716, \quad z_e = -17.001, \quad 
\zch_e = 17.716 - 17.001 = 0.715\MeV^{1/2}\,.
\end{equation}
A 96\% cancellation. The mass $m_e = 0.715^2 = 0.511\MeV$ is the 
square of something already small.

The near-vanishing of $m_e$ is a structural property of the Koide 
embedding, not fine-tuning: the energy balance allows perturbations 
as large as $z_0$. Every Koide triple has one member with 
$|\zch|/v_0$ small:

\begin{center}
\begin{tabular}{lccl}
\toprule
Triple & Lightest & $|\zch|/v_0$ & Character \\
\midrule
$(e, \mu, \tau)$ & $e$ & 0.040 & Almost-seed \\
$(-\pi, D_s, B)$ & $\pi$ (neg.) & 0.337 & Full triple \\
$(c, b, t)$ & $c$ & 0.236 & Full triple \\
$(-s, c, b)$ & $s$ (neg.) & 0.320 & Full triple \\
\bottomrule
\end{tabular}
\end{center}

The electron is by far the most extreme case, consistent with the 
lepton triple being nearest to its seed.

%======================================================================
\section{Symmetry Restoration}
\label{sec:restoration}
%======================================================================

The broken spectrum reassembles as symmetry-breaking effects are 
removed:

\begin{enumerate}
\item \textbf{Turn off QED:} Electromagnetic splittings vanish 
($\pi^\pm = \pi^0$, $K^\pm = K^0$). D-term effects disappear.
\item \textbf{Set $m_u = m_d$:} Exact isospin.
\item \textbf{Set $m_s = m_d$:} $SU(3)_V$ restored, 
$\pi = K = \eta$. Koide splitting $v_1 \to 0$; lepton partners 
degenerate.
\item \textbf{Send $m_q \to 0$:} Chiral limit, all meson charges 
$\to 0$. If sBootstrap SUSY restores, all light lepton masses 
vanish too. This is the seed configuration: $(0, 0, 0)$.
\end{enumerate}

The seed triples are the ``penultimate'' stage: one member has 
reached zero, but the other two retain the $2 - \sqrt{3}$ ratio. 
Full triples are the final broken state. The lepton triple, 
trapped between seed and bloom, reveals the hierarchy: 
chiral breaking $\to$ seed $\to$ bloom $\to$ full spectrum.

The direction of the flow also matters: approaching the chiral 
limit, the pion charge rises from $\zch = -11.81$ toward zero 
from below. It crosses zero exactly in the chiral limit. This 
means the pion was ``born negative''---its charge sign is set by 
the direction of chiral symmetry breaking, and the seed 
$(0, \pi, D_s)$ is the configuration at the moment of crossing.

%======================================================================
\section{Nature Abhors Chirality}
\label{sec:chirality}
%======================================================================

A pattern across the Standard Model: every chiral structure is 
eliminated in the infrared.

\begin{center}
\begin{tabular}{lcl}
\toprule
UV structure & IR survivor & Mechanism \\
\midrule
$SU(2)_L \times U(1)_Y$ & $U(1)_{\text{em}}$ & Higgs \\
$SU(N_f)_L \times SU(N_f)_R$ & $SU(N_f)_V$ & $\qqbar$ condensate \\
$U(1)_A$ & nothing & Anomaly \\
$\bar{\theta}$-vacuum & $|\bar{\theta}| < 10^{-10}$ & Unknown \\
Chiral color? & $SU(3)_C$ & Axigluon mass \\
\bottomrule
\end{tabular}
\end{center}

The sBootstrap adds: superspace chirality (boson vs.\ fermion 
distinction within a multiplet) is broken by the same condensate 
that breaks flavor chirality, with $f_{\text{SUSY}} = f_\pi$. The 
energy-balance condition $\langle z_k^2\rangle = z_0^2$ may be 
related: the breaking is ``as large as the background 
allows''---maximal, not perturbative.

%======================================================================
\section{Open Problems}
\label{sec:open}
%======================================================================

\subsection{The diquark statistics problem}
The sBootstrap requires antisymmetric $\overline{\mathbf{10}}$ 
diquarks; Pauli statistics for spin-0, $L=0$, color $\bar{3}$ 
diquarks requires symmetric flavor. Extra dynamical input---$L=1$ 
plus spin-1, instantons, or the $SO(32)$ embedding---is needed.

\subsection{What determines the Koide phase?}
The energy balance constrains but does not uniquely fix $\delta$. 
The value $\delta_\ell \approx 2/9$ for leptons must come from 
dynamics. Is it a fixed point?

\subsection{Connecting the three tuples}
The light sector (seed $\to$ bloom, scale $\sim \Lambda_{\text{QCD}}$) 
and heavy sector ($c,b,t$, scale $\sim v_{\text{EW}}$) must connect. 
The charm and bottom quarks appear in both, providing the links.

\subsection{The baryon sector}
In Seiberg's SQCD with $N_f = N_c = 3$, baryons are part of the 
low-energy spectrum. Their SUSY partners are ``baryonic fermions'' 
with no known candidates among observed particles.

\subsection{The strong CP problem}
The principle ``Nature abhors chirality'' predicts $\bar{\theta} = 0$ 
but provides no mechanism.

\subsection{The kaon--muon puzzle}
The SUSY splitting $m_K^2 - m_\mu^2 \approx (482\MeV)^2$ is much 
larger than $f_K^2 \approx (110\MeV)^2$, showing that the breaking 
is strongly flavor-dependent.

%======================================================================
\section{Conclusions}
\label{sec:conclusions}
%======================================================================

Three principal results:

\paragraph{1.\ The Koide formula is energy equipartition.}
$Q = 2/3$ is $\langle z_k^2\rangle = z_0^2$: perturbation energies 
match the vacuum energy. The signed-charge framework resolves the 
electron mass puzzle (near-cancellation, not fine-tuning) and reveals 
that mesons satisfy Koide with $(-\pi, D_s, B)$ to $0.1\%$.

\paragraph{2.\ Seed triples set the stage for SUSY breaking.}
The quark seed $(0,s,c)$ and meson seed $(0,\pi,D_s)$ share the 
universal ratio $2 - \sqrt{3}$ and operate at $v_0^2 \sim 
230$--$350\MeV$, close to the lepton triple's $v_0^2 = 314\MeV$. 
The massless member is the necessary condition for F-term 
(O'Raifeartaigh) SUSY breaking. Upon breaking, seeds bloom into 
full triples $(-s,c,b)$ and $(-\pi,D_s,B)$. The lepton triple 
is the ``almost-seed'' that has barely bloomed.

\paragraph{3.\ Two breaking mechanisms plus a heavy-sector origin.}
F-term breaking (via the massless seed member) generates the light 
sector; the Volkov--Akulov relation $f_{\text{VA}} = f_\pi$ is its 
low-energy realization, not a separate mechanism. D-term breaking 
($U(1)_{\text{em}}$ FI parameter) provides electromagnetic 
splittings as a perturbative correction. The heavy triple 
$(c,b,t)$ stands apart: no seed, no meson shadow, scale far above 
$\Lambda_{\text{QCD}}$, possibly connected to electroweak breaking.

The diquark sector, required to complete the SUSY multiplet count, 
must fill the antisymmetric $\overline{\mathbf{10}}$ of $SU(5)_f$ 
despite the apparent conflict with Fermi statistics---a problem 
that may ultimately require the $SO(32)$ embedding for its 
resolution.

The organizing principle is that each sector of the spectrum begins 
with a massless member (a seed), which is the necessary condition 
for spontaneous SUSY breaking. The breaking lifts the zero, 
generates the full triple, and the energy-balance condition 
$\langle z_k^2\rangle = z_0^2$ is preserved as a fixed point 
of the dynamics. The entire charged fermion and pseudoscalar meson 
spectrum is then determined by a small number of parameters: the 
vacuum charges $v_0$ of the three tuples, the Koide phases 
$\delta$ that control the bloom, and the SUSY-breaking scale 
$f_\pi$ that separates bosons from fermions.

\begin{thebibliography}{99}
\bibitem{Rivero2005}
A.~Rivero, ``The sBootstrap,'' viXra:0503.001 (2005);
``A prediction from the sBootstrap,'' viXra:0809.001 (2008);
``sBootstrap and SO(32),'' unpublished notes.

\bibitem{Rivero2012}
A.~Rivero, ``A new Koide tuple: strange, charm, bottom,'' 
arXiv:1111.7232 [hep-ph] (2011).

\bibitem{Koide1983}
Y.~Koide, ``New viewpoint of quark and lepton mass hierarchy,'' 
Phys.\ Rev.\ D \textbf{28}, 252 (1983).

\bibitem{CattoGursey}
S.~Catto and F.~G\"ursey, ``Algebraic treatment of effective 
supersymmetry,'' Nuovo Cim.\ A \textbf{86}, 201 (1985).

\bibitem{Miyazawa}
H.~Miyazawa, ``Baryon Number Changing Currents,'' 
Prog.\ Theor.\ Phys.\ \textbf{36}, 1266 (1966).

\bibitem{Harari}
H.~Harari, H.~Haut, J.~Weyers, ``Quark masses and Cabibbo 
angles,'' Phys.\ Lett.\ B \textbf{78}, 459 (1978).

\bibitem{Seiberg1995}
N.~Seiberg, ``Electric--magnetic duality in supersymmetric 
non-abelian gauge theories,'' Nucl.\ Phys.\ B \textbf{435}, 129 
(1995).

\bibitem{ORaifeartaigh}
L.~O'Raifeartaigh, ``Spontaneous symmetry breaking for chiral 
scalar superfields,'' Nucl.\ Phys.\ B \textbf{96}, 331 (1975).

\bibitem{FayetIliopoulos}
P.~Fayet and J.~Iliopoulos, ``Spontaneously broken supergauge 
symmetries and Goldstone spinors,'' Phys.\ Lett.\ B \textbf{51}, 
461 (1974).

\bibitem{Witten1983}
E.~Witten, ``Some inequalities among hadron masses,'' 
Phys.\ Rev.\ Lett.\ \textbf{51}, 2351 (1983).

\bibitem{VolkovAkulov1973}
D.~V.~Volkov and V.~P.~Akulov, ``Is the neutrino a Goldstone 
particle?,'' Phys.\ Lett.\ B \textbf{46}, 109 (1973).

\bibitem{GOR1968}
M.~Gell-Mann, R.~J.~Oakes, B.~Renner, ``Behavior of current 
divergences under $SU(3) \times SU(3)$,'' 
Phys.\ Rev.\ \textbf{175}, 2195 (1968).

\bibitem{ISS2006}
K.~Intriligator, N.~Seiberg, D.~Shih, ``Dynamical SUSY breaking 
in meta-stable vacua,'' JHEP \textbf{0604}, 021 (2006),
arXiv:hep-ph/0602239.

\bibitem{DGMLY1967}
T.~Das, G.~S.~Guralnik, V.~S.~Mathur, F.~E.~Low, J.~E.~Young,
``Electromagnetic mass difference of pions,''
Phys.\ Rev.\ Lett.\ \textbf{18}, 759 (1967).

\bibitem{FramptonGlashow1987}
P.~H.~Frampton and S.~L.~Glashow, ``Chiral color: an alternative 
to the standard model,'' Phys.\ Lett.\ B \textbf{190}, 157 (1987).

\end{thebibliography}

\end{document}
